\documentclass[12pt,a4paper]{article}

\usepackage[a4paper,margin=2.3cm,top=2.2cm,bottom=2.5cm]{geometry}
\usepackage{fontspec}
\usepackage{polyglossia}
\setmainlanguage{russian}
\setotherlanguage{english}
\setmainfont{DejaVu Serif}
\setsansfont{DejaVu Sans}
\setmonofont{DejaVu Sans Mono}

\usepackage{amsmath,amssymb,amsthm,mathtools,bm}
\usepackage{array,booktabs,longtable,tabularx,multirow}
\usepackage{enumitem}
\usepackage{graphicx}
\usepackage{tikz}
\usepackage{tcolorbox}
\tcbuselibrary{skins,breakable}
\usepackage{xcolor}
\usepackage{hyperref}
\usepackage{fancyhdr}
\usepackage{titlesec}
\usepackage{setspace}
\usepackage{microtype}
\usepackage{csquotes}

\hypersetup{
  colorlinks=true,
  linkcolor=blue!50!black,
  urlcolor=blue!50!black,
  citecolor=blue!50!black,
  pdftitle={Конспект по гладким многообразиям и римановой оптимизации},
  pdfauthor={OpenAI}
}

\definecolor{Accent}{HTML}{1F4E79}
\definecolor{AccentSoft}{HTML}{EEF4FA}
\definecolor{AccentMid}{HTML}{D8E6F3}
\definecolor{WarmGray}{HTML}{5F6368}
\definecolor{BoxGreen}{HTML}{EDF7F0}
\definecolor{BoxYellow}{HTML}{FFF7E8}
\definecolor{BoxRed}{HTML}{FFF0F0}

\titleformat{\section}
  {\Large\bfseries\color{Accent}}{\thesection.}{0.6em}{}
\titleformat{\subsection}
  {\large\bfseries\color{Accent}}{\thesubsection.}{0.6em}{}
\titleformat{\subsubsection}
  {\normalsize\bfseries\color{Accent}}{\thesubsubsection.}{0.6em}{}

\setlist[itemize]{leftmargin=1.6em,itemsep=0.25em,topsep=0.35em}
\setlist[enumerate]{leftmargin=1.8em,itemsep=0.25em,topsep=0.35em}
\setlength{\parindent}{0pt}
\setlength{\parskip}{0.55em}
\onehalfspacing
\setlength{\headheight}{25pt}

\pagestyle{fancy}
\fancyhf{}
\fancyhead[L]{\textcolor{WarmGray}{Гладкие многообразия и риманова оптимизация}}
\fancyhead[R]{\textcolor{WarmGray}{\thepage}}
\renewcommand{\headrulewidth}{0.3pt}
\renewcommand{\headrule}{\hbox to\headwidth{\color{AccentMid}\leaders\hrule height \headrulewidth\hfill}}

\newtheoremstyle{ruplain}
  {0.7em}{0.7em}{\itshape}{} {\bfseries\color{Accent}}{.}{0.4em}{}
\newtheoremstyle{rudef}
  {0.7em}{0.7em}{}{} {\bfseries\color{Accent}}{.}{0.4em}{}
\theoremstyle{rudef}
\newtheorem{definition}{Определение}[section]
\newtheorem{remark}[definition]{Замечание}
\newtheorem{example}[definition]{Пример}
\newtheorem{proposition}[definition]{Утверждение}
\theoremstyle{ruplain}
\newtheorem{theorem}[definition]{Теорема}
\newtheorem{lemma}[definition]{Лемма}

\newtcolorbox{summarybox}{
  breakable,
  colback=AccentSoft,
  colframe=AccentMid,
  boxrule=0.6pt,
  arc=2mm,
  left=1.4mm,right=1.4mm,top=1.2mm,bottom=1.2mm
}

\newtcolorbox{ideabox}{
  breakable,
  colback=BoxGreen,
  colframe=green!35!black,
  title=Идея,
  fonttitle=\bfseries,
  boxrule=0.6pt,
  arc=2mm,
  left=1.4mm,right=1.4mm,top=1.2mm,bottom=1.2mm
}

\newtcolorbox{warningbox}{
  breakable,
  colback=BoxYellow,
  colframe=orange!55!black,
  title=Важно,
  fonttitle=\bfseries,
  boxrule=0.6pt,
  arc=2mm,
  left=1.4mm,right=1.4mm,top=1.2mm,bottom=1.2mm
}

\newtcolorbox{formulaibox}{
  breakable,
  colback=white,
  colframe=AccentMid,
  title=Ключевая формула,
  fonttitle=\bfseries,
  boxrule=0.6pt,
  arc=2mm,
  left=1.4mm,right=1.4mm,top=1.2mm,bottom=1.2mm
}

\newcommand{\R}{\mathbb{R}}
\newcommand{\grad}{\operatorname{grad}}
\newcommand{\Sym}{\operatorname{Sym}}
\newcommand{\rank}{\operatorname{rank}}
\newcommand{\tr}{\operatorname{tr}}
\newcommand{\Span}{\operatorname{span}}
\newcommand{\Proj}{\operatorname{Proj}}
\newcommand{\St}{\mathrm{St}}
\newcommand{\Gr}{\mathrm{Gr}}
\newcommand{\Exp}{\operatorname{Exp}}
\newcommand{\Retr}{\operatorname{Retr}}
\newcommand{\diag}{\operatorname{diag}}
\newcommand{\ip}[2]{\left\langle #1,#2\right\rangle}
\newcommand{\norm}[1]{\left\lVert #1\right\rVert}
\newcommand{\set}[1]{\left\{#1\right\}}
\newcommand{\parent}[1]{\left(#1\right)}
\newcommand{\brak}[1]{\left[#1\right]}

\begin{document}

\begin{titlepage}
\thispagestyle{empty}
\vspace*{2.8cm}

{\Large\color{Accent}\bfseries Конспект по гладким многообразиям\par}
\vspace{0.35cm}
{\Large\color{Accent}\bfseries и введение в риманову оптимизацию\par}

\vspace{1.2cm}
{\large Подробный разбор тем: гладкие и вложенные многообразия, касательные пространства, риманова метрика, риманов градиент, проекция на касательное пространство, ретракции, риманов градиентный спуск, многообразие матриц фиксированного ранга, Stiefel и Grassmann.\par}

\vspace{1.2cm}
\begin{summarybox}
\textbf{Как читать этот текст.} Этот конспект написан в жанре \emph{«восстановить картину целиком»}. Здесь не только даются определения, но и объясняется, зачем они нужны, как связаны между собой и как из геометрии естественно возникает риманова оптимизация.

Основой послужили загруженные лекции и семинары по классической дифференциальной геометрии (Иванов, Фоменко, Пенской). Однако темы, которых в этих материалах нет или которые там обсуждаются лишь косвенно, --- прежде всего риманов градиент, ретракции, градиентный спуск на многообразиях и матричные многообразия --- здесь разобраны отдельно и заметно подробнее.
\end{summarybox}

\vfill
{\large \textbf{Формат итогового текста:}\ учебниковый конспект с интуицией, формулами, примерами и короткими «шпаргалками» по ключевым объектам.\par}

\vspace{1.0cm}
{\large Stockholm / 2026\par}
\end{titlepage}

\tableofcontents
\newpage

\section{Зачем вообще нужны многообразия в задачах оптимизации}

Во многих задачах переменная не может быть произвольной точкой пространства $\R^n$. Она обязана удовлетворять \emph{гладким ограничениям}. Например:
\begin{itemize}
  \item единичный вектор $x \in S^{n-1}$, то есть $\norm{x}=1$;
  \item матрица с ортонормированными столбцами $X^TX=I$;
  \item подпространство фиксированной размерности;
  \item матрица фиксированного ранга $r$.
\end{itemize}

Если смотреть на такие ограничения геометрически, то множество допустимых точек обычно оказывается не «кусочком линейного пространства», а \emph{гладким многообразием}. На нём есть свои касательные направления, своя метрика, свои геодезические и, следовательно, свой «правильный» градиент.

\begin{ideabox}
Обычный евклидов градиент отвечает на вопрос: \emph{в какую сторону функция быстрее всего растёт в объёмлющем пространстве?}

Риманов градиент отвечает на другой вопрос: \emph{в какую допустимую касательную сторону функция быстрее всего растёт вдоль самого многообразия?}
\end{ideabox}

Поэтому риманова оптимизация --- это не какая-то «другая математика», а аккуратное приспособление привычных методов оптимизации к ситуации, когда допустимое множество криволинейно.

\section{Гладкие многообразия}

\subsection{Интуитивная картина}

Грубая идея проста: $m$-мерное гладкое многообразие --- это объект, который \emph{локально похож на $\R^m$}. То есть возле каждой точки можно ввести координаты, в которых наш объект выглядит как открытое множество в $\R^m$.

У сферы $S^2$ глобально нет одной хорошей карты без разрывов, но в окрестности каждой точки она выглядит как кусочек плоскости. У тора локально всё также похоже на плоскость. У окружности локально всё похоже на прямую.

Главная мысль: \textbf{размерность многообразия} --- это не размерность объёмлющего пространства, а число локальных координат, необходимых для описания точки на самом объекте.

\subsection{Карты и атласы}

\begin{definition}
Топологическое пространство $M$ называется $m$-мерным \emph{топологическим многообразием}, если для каждой точки $p\in M$ существует окрестность $U\subset M$ и гомеоморфизм
\[
\varphi:U\to \varphi(U)\subset \R^m,
\]
где $\varphi(U)$ --- открытое множество в $\R^m$.
\end{definition}

Пара $(U,\varphi)$ называется \emph{картой}. Если набор карт покрывает всё $M$, то он называется \emph{атласом}.

Чтобы на многообразии можно было дифференцировать, нужны гладкие переходы между координатами.

\begin{definition}
Атлас называется \emph{гладким}, если для любых двух карт $(U,\varphi)$ и $(V,\psi)$ с непустым пересечением отображение перехода
\[
\psi\circ \varphi^{-1}:\varphi(U\cap V)\to \psi(U\cap V)
\]
является гладким. Многообразие с таким атласом называется \emph{гладким многообразием}.
\end{definition}

\subsection{Почему координаты здесь не просто украшение}

В линейном пространстве у нас есть глобальные координаты. На многообразии одной глобальной системы координат, вообще говоря, нет. Поэтому локальные карты здесь --- не дополнительный комфорт, а \emph{основной способ определения гладкости}.

Именно через карты определяются:
\begin{itemize}
  \item гладкие функции на многообразии;
  \item гладкие отображения между многообразиями;
  \item касательные векторы;
  \item компоненты римановой метрики;
  \item координатные формулы для градиента, связности и т.д.
\end{itemize}

\subsection{Стандартные примеры}

\begin{example}[Сфера]
Сфера
\[
S^{n-1}=\set{x\in \R^n: \norm{x}=1}
\]
является $(n-1)$-мерным гладким многообразием. Локальные координаты можно задавать, например, через стереографическую проекцию.
\end{example}

\begin{example}[Тор]
Стандартный тор в $\R^3$ --- двумерное гладкое многообразие. Его можно параметризовать двумя углами.
\end{example}

\begin{example}[Пространство обратимых матриц]
Множество
\[
GL(n)=\set{A\in \R^{n\times n}: \det A\neq 0}
\]
есть открытое подмножество пространства матриц $\R^{n\times n}$, а потому автоматически гладкое многообразие размерности $n^2$.
\end{example}

\begin{example}[Stiefel]
\[
\St(p,n)=\set{X\in \R^{n\times p}: X^TX=I_p}
\]
--- многообразие матриц с ортонормированными столбцами.
\end{example}

\begin{example}[Grassmann]
$\Gr(p,n)$ --- множество всех $p$-мерных подпространств пространства $\R^n$.
\end{example}

\begin{example}[Матрицы фиксированного ранга]
\[
\mathcal M_r=\set{X\in \R^{m\times n}: \rank X=r}
\]
--- гладкое многообразие. Именно оно играет ключевую роль в низкоранговой оптимизации.
\end{example}

\section{Вложенные многообразия и подмногообразия в $\R^n$}

Именно этот случай чаще всего нужен в оптимизации: допустимое множество сидит внутри обычного евклидова пространства, и хочется использовать структуру объёмлющего пространства.

\subsection{Три локально эквивалентные точки зрения}

В лекциях по Иванову подчёркивается локальная эквивалентность трёх способов задания кривых и поверхностей: как образа параметризации, как графика и как множества решений невырожденной системы уравнений. Эта же идея переносится на многообразия более общей размерности.

Почти всегда подмногообразие в $\R^n$ можно локально понимать тремя способами:
\begin{enumerate}
  \item как образ гладкого отображения $r:U\subset \R^m\to \R^n$ ранга $m$;
  \item как график гладкой функции;
  \item как множество решений системы уравнений $F(x)=0$ при условии максимального ранга $dF$.
\end{enumerate}

\subsection{Параметрическое задание}

\begin{definition}
Пусть $U\subset \R^m$ открыто, а $r:U\to \R^n$ --- гладкое отображение. Если дифференциал $dr_u$ имеет ранг $m$ во всех точках $u\in U$ и отображение $r$ задаёт гомеоморфизм $U$ со своим образом, то образ $r(U)$ называется \emph{вложенным $m$-мерным подмногообразием} в $\R^n$.
\end{definition}

Интуитивно это значит, что параметры $u^1,\dots,u^m$ действительно описывают $m$ независимых направлений вдоль поверхности.

\subsection{Задание уравнениями и теорема о неявной функции}

Пусть дано отображение
\[
F:\R^n\to \R^k,
\qquad
M=\set{x\in \R^n:F(x)=0}.
\]
Если в каждой точке $x\in M$ матрица Якоби $dF_x$ имеет ранг $k$, то локально множество решений можно разрешить по $k$ переменным. Поэтому $M$ локально выглядит как график функции от $n-k$ переменных.

\begin{theorem}[Критерий гладкого подмногообразия по уравнениям]
Если $F:\R^n\to \R^k$ гладко и $\rank dF_x=k$ для всех $x\in M=F^{-1}(0)$, то $M$ является гладким подмногообразием размерности $n-k$.
\end{theorem}

\begin{example}[Сфера]
Сфера задаётся одним уравнением
\[
F(x)=x^Tx-1=0.
\]
Здесь $dF_x(h)=2x^Th$, и при $x\in S^{n-1}$ этот функционал не нулевой. Следовательно, $S^{n-1}$ --- гладкое подмногообразие коразмерности $1$.
\end{example}

\subsection{Погружение и вложение}

По материалам Иванова полезно различать два понятия.

\begin{definition}
Гладкое отображение $F:M\to N$ называется \emph{погружением} (иммерсией), если $dF_p:T_pM\to T_{F(p)}N$ инъективно для всех $p\in M$.
\end{definition}

\begin{definition}
Погружение называется \emph{вложением}, если оно дополнительно задаёт гомеоморфизм $M$ на образ $F(M)$ с индуцированной топологией.
\end{definition}

Главная разница: погружение контролирует \emph{дифференциальную} структуру, а вложение ещё и \emph{топологически} запрещает ненужные самосклейки.

\subsection{Практический критерий: касательные пространства как ядро ограничений}

Если
\[
M=\set{x\in \R^n:F(x)=0}, \qquad F=(F_1,\dots,F_k),
\]
и $\rank dF_x=k$, то любой допустимый касательный вектор $\xi$ должен удовлетворять линейной системе
\[
dF_x[\xi]=0.
\]
Значит,
\[
T_xM=\ker dF_x.
\]

Это одна из самых важных формул во всём прикладном сюжете.

\begin{example}[Сфера]
Для $S^{n-1}$ имеем $F(x)=x^Tx-1$. Тогда
\[
dF_x[\xi]=2x^T\xi,
\]
поэтому
\[
T_xS^{n-1}=\set{\xi\in \R^n: x^T\xi=0}.
\]
То есть касательное пространство к сфере --- это гиперплоскость, ортогональная радиусу.
\end{example}

\section{Касательное пространство}

\subsection{Через скорости кривых}

В классической дифференциальной геометрии для подмногообразий в $\R^n$ естественно начать с простейшего определения.

\begin{definition}
Пусть $M\subset \R^n$ --- гладкое подмногообразие и $x\in M$. Касательным вектором в точке $x$ называется вектор скорости $\dot \gamma(0)$ некоторой гладкой кривой $\gamma:(-\varepsilon,\varepsilon)\to M$ такой, что $\gamma(0)=x$. Множество всех таких векторов называется \emph{касательным пространством} и обозначается $T_xM$.
\end{definition}

Для вложенных многообразий это определение очень наглядно: допустимая кривая должна оставаться на многообразии, а её мгновенная скорость и есть касательное направление.

\subsection{Через координаты}

Если $x=r(u)$ и $r:U\subset \R^m\to M\subset \R^n$ --- локальная параметризация, то производные
\[
\frac{\partial r}{\partial u^1}(u),\dots,\frac{\partial r}{\partial u^m}(u)
\]
образуют базис касательного пространства. Поэтому любой касательный вектор имеет вид
\[
\xi=\sum_{i=1}^m \xi^i\frac{\partial r}{\partial u^i}(u).
\]

Здесь числа $\xi^1,\dots,\xi^m$ --- это координаты касательного вектора в базисе, связанном с выбранной картой.

\subsection{Через производные по функциям: абстрактная точка зрения}

Для абстрактного многообразия определение через векторы в $\R^n$ уже не работает, потому что само многообразие никуда не обязано быть вложено. Тогда касательный вектор можно определить как дифференцирование функций.

\begin{definition}
Касательным вектором $v\in T_xM$ можно считать линейный функционал на пространстве гладких функций $f\in C^{\infty}(M)$, удовлетворяющий правилу Лейбница:
\[
v(fg)=v(f)\,g(x)+f(x)\,v(g).
\]
Такие функционалы называют \emph{деривациями} в точке $x$.
\end{definition}

Эта точка зрения особенно удобна, когда нужно определять дифференциал отображения и тензорные объекты.

\subsection{Почему все определения эквивалентны}

Схематически картина такова:
\begin{itemize}
  \item скорость кривой определяет производную функций вдоль этой кривой;
  \item в локальных координатах всякий касательный вектор записывается как линейная комбинация $\partial/\partial x^i$;
  \item правило преобразования координат обеспечивает независимость от выбора карты.
\end{itemize}

Именно поэтому можно свободно переключаться между геометрической, координатной и алгебраической точками зрения.

\subsection{Быстрые способы находить $T_xM$ на практике}

Для задач оптимизации и вычислений обычно используются три шаблона.

\begin{summarybox}
\textbf{Шаблон 1: через параметризацию.}
Если $x=r(u)$, то
\[
T_xM = \Span\set{\frac{\partial r}{\partial u^1},\dots,\frac{\partial r}{\partial u^m}}.
\]

\textbf{Шаблон 2: через ограничения.}
Если $M=\set{x:F(x)=0}$ и $dF_x$ имеет полный ранг, то
\[
T_xM=\ker dF_x.
\]

\textbf{Шаблон 3: через ортогональность нормалям.}
Если ограничения заданы скалярными функциями $F_1,\dots,F_k$, то
\[
T_xM = \set{\xi:\ \ip{\nabla F_i(x)}{\xi}=0,\ i=1,\dots,k}.
\]
\end{summarybox}

\subsection{Примеры}

\begin{example}[Сфера $S^{n-1}$]
\[
T_xS^{n-1}=\set{\xi\in \R^n:\ x^T\xi=0}.
\]
\end{example}

\begin{example}[Stiefel]
Многообразие задаётся матричным ограничением $X^TX=I_p$. Дифференцируя его в направлении $\Xi$, получаем
\[
X^T\Xi+\Xi^TX=0.
\]
Следовательно,
\[
T_X\St(p,n)=\set{\Xi\in \R^{n\times p}: X^T\Xi+\Xi^TX=0}.
\]
\end{example}

\begin{example}[Grassmann в ортонормированном представителе]
Если точка задаётся ортонормированной матрицей $X\in \St(p,n)$, то касательные направления к самому пространству подпространств можно выбирать горизонтально:
\[
T_{[X]}\Gr(p,n)\cong \set{\Xi\in \R^{n\times p}: X^T\Xi=0}.
\]
\end{example}

\section{Дифференциал гладкого отображения}

Пусть $F:M\to N$ --- гладкое отображение между гладкими многообразиями.

\begin{definition}
\emph{Дифференциал} отображения $F$ в точке $x\in M$ --- это линейное отображение
\[
dF_x:T_xM\to T_{F(x)}N,
\]
которое переводит касательный вектор в скорость образа соответствующей кривой:
\[
dF_x(\dot \gamma(0)) = \frac{d}{dt}\Big|_{t=0} F(\gamma(t)).
\]
\end{definition}

В координатах дифференциал задаётся матрицей Якоби, а для функций $f:M\to \R$ он становится ковектором
\[
df_x:T_xM\to \R.
\]

\begin{warningbox}
Очень важно различать:
\begin{itemize}
  \item \textbf{дифференциал функции} $df_x$ --- линейный функционал на $T_xM$;
  \item \textbf{градиент} $\grad f(x)$ --- вектор в $T_xM$.
\end{itemize}
Они связаны римановой метрикой, но это не одно и то же.
\end{warningbox}

\section{Риманова метрика}

\subsection{Что она делает геометрически}

На каждом касательном пространстве $T_xM$ риманова метрика задаёт скалярное произведение. Значит, в каждой точке можно измерять:
\begin{itemize}
  \item длины касательных векторов;
  \item углы между касательными векторами;
  \item длины кривых на многообразии;
  \item норму градиента;
  \item направление наискорейшего роста/убывания вдоль многообразия.
\end{itemize}

\begin{definition}
Римановой метрикой на гладком многообразии $M$ называется семейство скалярных произведений
\[
g_x(\cdot,\cdot):T_xM\times T_xM\to \R,
\qquad x\in M,
\]
гладко зависящее от точки $x$.
\end{definition}

\subsection{Координатная запись}

Если в карте $(x^1,\dots,x^m)$ взять координатные векторы
\[
\frac{\partial}{\partial x^1},\dots,\frac{\partial}{\partial x^m},
\]
то компоненты метрики определяются формулами
\[
g_{ij}(x)=g_x\!\left(\frac{\partial}{\partial x^i},\frac{\partial}{\partial x^j}\right).
\]
Тогда для касательных векторов
\[
\xi=\sum_i \xi^i \frac{\partial}{\partial x^i},
\qquad
\eta=\sum_j \eta^j \frac{\partial}{\partial x^j}
\]
имеем
\[
g_x(\xi,\eta)=\sum_{i,j} g_{ij}(x)\,\xi^i\eta^j.
\]

Матрица $G(x)=(g_{ij}(x))$ симметрична и положительно определена.

\subsection{Длина кривой и угол}

Если $\gamma:[a,b]\to M$ --- гладкая кривая, то её длина определяется как
\[
L(\gamma)=\int_a^b \sqrt{g_{\gamma(t)}(\dot\gamma(t),\dot\gamma(t))}\,dt.
\]
Угол между двумя ненулевыми касательными векторами $\xi,\eta\in T_xM$ задаётся формулой
\[
\cos\theta = \frac{g_x(\xi,\eta)}{\norm{\xi}_x\norm{\eta}_x},
\qquad
\norm{\xi}_x=\sqrt{g_x(\xi,\xi)}.
\]

\subsection{Индуцированная метрика}

Если $M\subset \R^n$ вложено в евклидово пространство, то самый естественный выбор --- взять метрику, индуцированную стандартным скалярным произведением в $\R^n$:
\[
g_x(\xi,\eta)=\ip{\xi}{\eta}_{\R^n},
\qquad \xi,\eta\in T_xM.
\]

Это особенно удобно для оптимизации, потому что позволяет использовать ambient-вычисления: сначала всё посчитать в $\R^n$, а затем спроецировать в касательное пространство.

\subsection{Индуцированная метрика в параметризации}

Если $x=r(u^1,\dots,u^m)$ --- локальная параметризация вложенного многообразия, то
\[
g_{ij}(u)=\ip{\frac{\partial r}{\partial u^i},\frac{\partial r}{\partial u^j}}.
\]

Это и есть обобщение первой фундаментальной формы поверхности на многомерный случай.

\section{Евклидов градиент и риманов градиент}

\subsection{Дифференциал функции как исходная величина}

Пусть $f:M\to \R$ --- гладкая функция. В точке $x\in M$ её дифференциал $df_x$ --- это линейный функционал на $T_xM$. Если $\xi\in T_xM$, то число $df_x[\xi]$ --- это производная функции в направлении $\xi$.

Сами по себе производные по направлениям образуют \emph{ковектор}. Чтобы превратить их в \emph{вектор}, нужна метрика.

\subsection{Евклидов случай}

В $\R^n$ со стандартным скалярным произведением градиент $\nabla f(x)$ определяется условием
\[
df_x[h]=\ip{\nabla f(x)}{h}
\qquad \text{для всех } h\in \R^n.
\]

Если $f=f(x^1,\dots,x^n)$, то
\[
\nabla f = \begin{pmatrix} \partial_1 f \\ \vdots \\ \partial_n f \end{pmatrix}.
\]

Отрицательный градиент $-\nabla f$ есть направление наискорейшего убывания функции в евклидовой норме.

\subsection{Риманов градиент: определение}

\begin{definition}
Римановым градиентом функции $f:M\to \R$ в точке $x$ называется единственный вектор $\grad f(x)\in T_xM$, удовлетворяющий условию
\[
df_x[\xi]=g_x(\grad f(x),\xi)
\qquad \text{для всех } \xi\in T_xM.
\]
\end{definition}

Это прямой аналог евклидовой формулы, только скалярное произведение теперь берётся не в $\R^n$, а в касательном пространстве с метрикой $g_x$.

\begin{ideabox}
Дифференциал $df_x$ --- это «список производных по всем допустимым направлениям». Градиент $\grad f(x)$ --- это вектор, который кодирует тот же объект после отождествления касательного пространства с его двойственным при помощи метрики.
\end{ideabox}

\subsection{Координатная формула}

Если в локальных координатах
\[
df = \sum_{i=1}^m \frac{\partial f}{\partial x^i}\,dx^i,
\]
а матрица метрики равна $G=(g_{ij})$, то для координат риманова градиента
\[
\grad f = \sum_{i=1}^m (\grad f)^i \frac{\partial}{\partial x^i}
\]
получаем систему
\[
\sum_j g_{ij}(\grad f)^j = \frac{\partial f}{\partial x^i}.
\]
Значит,
\[
(\grad f)^i = \sum_j g^{ij}\frac{\partial f}{\partial x^j},
\]
где $(g^{ij})=G^{-1}$.

\begin{formulaibox}
В координатах
\[
\boxed{\ \grad f = \sum_{i,j} g^{ij}\,\partial_j f\,\frac{\partial}{\partial x^i}.\ }
\]
В евклидовых координатах $g_{ij}=\delta_{ij}$, и эта формула превращается в обычный градиент.
\end{formulaibox}

\subsection{Вложенное многообразие: связь с евклидовым градиентом}

Пусть $M\subset \R^n$ --- вложенное многообразие с индуцированной метрикой, а $f:M\to \R$. Возьмём гладкое продолжение $\bar f$ функции $f$ в окрестность точки $x$ в $\R^n$. Тогда
\[
df_x[\xi]=d\bar f_x[\xi]=\ip{\nabla \bar f(x)}{\xi}
\qquad \text{для всех } \xi\in T_xM.
\]

С другой стороны,
\[
df_x[\xi]=g_x(\grad f(x),\xi)=\ip{\grad f(x)}{\xi}.
\]

Следовательно, риманов градиент --- это просто ортогональная проекция евклидова градиента на касательное пространство:
\[
\boxed{\ \grad f(x)=\Proj_{T_xM}\parent{\nabla \bar f(x)}.\ }
\]

Эта формула --- центральная для embedded-оптимизации.

\subsection{Почему именно этот вектор является «правильным»}

Если фиксировать норму, заданную метрикой $g_x$, то среди всех единичных касательных векторов $\xi$ производная $df_x[\xi]$ максимальна именно на направлении
\[
\xi = \frac{\grad f(x)}{\norm{\grad f(x)}_x},
\]
а минимальна на противоположном направлении. Поэтому в римановой оптимизации отрицательный риманов градиент есть естественное направление спуска.

\section{Проекция на касательное пространство}

\subsection{Общий принцип}

Пусть $M\subset \R^n$ вложено в евклидово пространство. Если уже найден какой-то вектор $z\in \R^n$, то его риманов аналог вдоль многообразия обычно находится как
\[
\Proj_{T_xM}(z).
\]

В задачах оптимизации этим $z$ чаще всего является евклидов градиент.

\subsection{Сфера}

Для единичной сферы $S^{n-1}$
\[
T_xS^{n-1}=\set{\xi: x^T\xi=0}.
\]
Разложим $z\in \R^n$ на тангенциальную и нормальную части:
\[
z = \underbrace{\parent{z-(x^Tz)x}}_{\text{касательная часть}} + \underbrace{(x^Tz)x}_{\text{нормальная часть}}.
\]
Следовательно,
\[
\boxed{\ \Proj_{T_xS^{n-1}}(z)=z-(x^Tz)x.\ }
\]

\subsection{Stiefel}

Напомним, что
\[
\St(p,n)=\set{X\in \R^{n\times p}: X^TX=I_p},
\qquad
T_X\St(p,n)=\set{\Xi: X^T\Xi+\Xi^TX=0}.
\]

Ортогональная проекция матрицы $Z\in \R^{n\times p}$ на $T_X\St(p,n)$ при стандартной метрике Фробениуса имеет вид
\[
\boxed{\ \Proj_{T_X\St}(Z)= Z - X\Sym(X^TZ),\qquad \Sym(A)=\frac{A+A^T}{2}.\ }
\]

Проверка проста: после подстановки условие касательности выполняется автоматически.

\subsection{Grassmann}

Точку $\Gr(p,n)$ удобно представлять ортонормированной матрицей $X\in \St(p,n)$, столбцы которой порождают нужное подпространство. Касательные направления можно брать из пространства
\[
\set{\Xi\in \R^{n\times p}: X^T\Xi=0}.
\]
Тогда проекция имеет особенно простой вид:
\[
\boxed{\ \Proj_{T_{[X]}\Gr}(Z) = (I-XX^T)Z.\ }
\]

\subsection{Многообразие матриц фиксированного ранга}

Пусть $X\in \mathcal M_r$ и его компактное SVD имеет вид
\[
X=U\Sigma V^T,
\qquad U\in \St(r,m),\ V\in \St(r,n),\ \Sigma\in \R^{r\times r}.
\]
Обозначим ортогональные проекторы на столбцовые пространства через
\[
P_U=UU^T,\qquad P_V=VV^T.
\]
Тогда ортогональная проекция $Z\in \R^{m\times n}$ на касательное пространство в точке $X$ равна
\[
\boxed{\ \Proj_{T_X\mathcal M_r}(Z)=P_UZ + ZP_V - P_U Z P_V.\ }
\]

Именно эта формула даёт риманов градиент для низкоранговых задач.

\section{Геодезические, экспоненциальное отображение и ретракции}

\subsection{Зачем нужен ещё один шаг после проекции}

Даже если мы нашли хороший касательный вектор $\xi\in T_xM$, точка
\[
x-\alpha \xi
\]
обычно уже \emph{не лежит} на многообразии. Поэтому после шага в касательном пространстве нужно вернуть точку обратно на $M$.

Самый «геометрически честный» способ --- идти по геодезической.

\subsection{Геодезическая и экспоненциальное отображение}

Геодезическую можно понимать как кривую, которая локально идёт «прямее всего» с точки зрения метрики. В евклидовом пространстве геодезические --- обычные прямые. На сфере --- большие окружности.

Если $\gamma_{x,\xi}(t)$ --- геодезическая, стартующая из точки $x$ со скоростью $\xi$, то определяют экспоненциальное отображение
\[
\Exp_x(\xi)=\gamma_{x,\xi}(1).
\]

То есть экспонента берёт касательный вектор и говорит, куда мы придём через единицу времени, если пойдём из $x$ по истинной геодезической с начальной скоростью $\xi$.

\begin{example}[Евклидово пространство]
В $\R^n$ геодезическая имеет вид $\gamma(t)=x+t\xi$, поэтому
\[
\Exp_x(\xi)=x+\xi.
\]
Здесь экспоненциальное отображение совпадает с обычным прибавлением вектора.
\end{example}

\begin{example}[Сфера]
На единичной сфере при $\xi\in T_xS^{n-1}$, $\xi\neq 0$,
\[
\Exp_x(\xi)= \cos\norm{\xi}\,x + \sin\norm{\xi}\,\frac{\xi}{\norm{\xi}}.
\]
\end{example}

\subsection{Почему экспонента неудобна вычислительно}

Экспонента концептуально прекрасна, но на практике часто дорога:
\begin{itemize}
  \item её формула может быть сложной;
  \item нужно решать геодезическое уравнение;
  \item для матричных многообразий точная экспонента часто существенно тяжелее простого шага QR или SVD.
\end{itemize}

Отсюда возникает идея заменить экспоненту \emph{локально эквивалентным} более дешёвым отображением.

\subsection{Ретракция}

\begin{definition}
Отображение
\[
\Retr_x:T_xM\to M
\]
называется \emph{ретракцией}, если выполняются два условия:
\[
\Retr_x(0_x)=x,
\qquad
D\Retr_x(0_x)=\mathrm{id}_{T_xM}.
\]
То есть ретракция в нуле совпадает с точкой $x$ и имеет тот же первый дифференциал, что и экспонента.
\end{definition}

На уровне интуиции: ретракция --- это \emph{первого порядка точная аппроксимация} истинного геодезического шага.

\begin{warningbox}
Ретракция не обязана быть геодезической. Её задача скромнее: корректно вернуть точку на многообразие и правильно совпасть с экспонентой в первом порядке по малому касательному вектору.
\end{warningbox}

\subsection{Примеры ретракций}

\begin{example}[Сфера]
Для $x\in S^{n-1}$ и $\xi\in T_xS^{n-1}$
\[
\Retr_x(\xi)=\frac{x+\xi}{\norm{x+\xi}}.
\]
Это самый естественный способ: сделать евклидов шаг и затем нормировать.
\end{example}

\begin{example}[Stiefel: QR-ретракция]
Если $X\in \St(p,n)$ и $\Xi\in T_X\St(p,n)$, можно положить
\[
\Retr_X(\Xi)=\operatorname{qf}(X+\Xi),
\]
где $\operatorname{qf}$ обозначает $Q$-фактор из QR-разложения.
\end{example}

\begin{example}[Stiefel: полярная ретракция]
\[
\Retr_X(\Xi)=(X+\Xi)\bigl(I_p+\Xi^T\Xi\bigr)^{-1/2}.
\]
Она обычно лучше согласована с геометрией, чем QR, но немного дороже.
\end{example}

\begin{example}[Grassmann]
Для представителя $X\in \St(p,n)$ можно использовать ту же идею:
\[
\Retr_X(\Xi)=\operatorname{qf}(X+\Xi),
\]
а затем понимать результат как новое подпространство.
\end{example}

\begin{example}[Фиксированный ранг]
Для $X\in \mathcal M_r$ и $\Xi\in T_X\mathcal M_r$ можно взять ретракцию
\[
\Retr_X(\Xi)=\Pi_r(X+\Xi),
\]
где $\Pi_r$ --- лучшая ранга-$r$ аппроксимация (усечённое SVD). Это очень важная практическая ретракция.
\end{example}

\section{Идея риманова градиентного спуска}

\subsection{Постановка задачи}

Пусть нужно минимизировать функцию
\[
f:M\to \R
\]
на римановом многообразии $(M,g)$.

В евклидовом пространстве обычный градиентный спуск имеет вид
\[
x_{k+1}=x_k-\alpha_k\nabla f(x_k).
\]

На многообразии эта запись бессмысленна буквально, потому что сумма точки и касательного вектора вообще не обязана принадлежать $M$. Правильная схема состоит из двух шагов:
\begin{enumerate}
  \item вычислить риманов градиент $\grad f(x_k)\in T_{x_k}M$;
  \item вернуться на многообразие с помощью ретракции.
\end{enumerate}

\subsection{Один шаг метода}

\begin{formulaibox}
Один шаг риманова градиентного спуска имеет вид
\[
\boxed{\ x_{k+1}=\Retr_{x_k}\!\bigl(-\alpha_k\,\grad f(x_k)\bigr).\ }
\]
\end{formulaibox}

Если метрика индуцирована из объёмлющего пространства, то часто удобно записывать это так:
\[
\grad f(x_k)=\Proj_{T_{x_k}M}(\nabla \bar f(x_k)),
\qquad
x_{k+1}=\Retr_{x_k}\!\bigl(-\alpha_k\grad f(x_k)\bigr).
\]

\subsection{Почему это действительно аналог обычного градиентного спуска}

Потому что структура полностью та же:
\begin{enumerate}
  \item сначала берём направление наискорейшего роста;
  \item меняем знак, чтобы получить направление спуска;
  \item выбираем длину шага $\alpha_k$;
  \item обновляем точку.
\end{enumerate}

Отличие только в том, что «разрешённые направления» теперь ограничены касательным пространством, а обновление точки должно учитывать кривизну допустимого множества.

\subsection{Псевдокод}

\begin{summarybox}
\textbf{Риманов градиентный спуск}
\begin{enumerate}
  \item выбрать начальную точку $x_0\in M$;
  \item для $k=0,1,2,\dots$ делать:
  \begin{enumerate}[label=\arabic*.]
    \item вычислить евклидов градиент $\nabla \bar f(x_k)$ или сразу $df_{x_k}$;
    \item получить риманов градиент $\grad f(x_k)$;
    \item выбрать шаг $\alpha_k$ (фиксированный, backtracking, Armijo и т.д.);
    \item сделать шаг
    \[
    x_{k+1}=\Retr_{x_k}\!\bigl(-\alpha_k\grad f(x_k)\bigr);
    \]
    \item остановиться, если $\norm{\grad f(x_k)}$ мало.
  \end{enumerate}
\end{enumerate}
\end{summarybox}

\subsection{Сравнение с проецированным градиентом}

Риманов градиентный спуск близок по духу к проецированному градиенту, но не совпадает с ним буквально.

Проецированный градиент в евклидовом пространстве часто делает шаг в ambient-пространстве, а потом проектирует \emph{точку} обратно на множество. Риманов подход сначала проектирует \emph{направление} в касательное пространство, а затем использует геометрически согласованную процедуру возвращения на многообразие.

При гладких подмногообразиях и хороших ретракциях эти идеи нередко оказываются очень близки, но риманова формулировка концептуально чище и лучше масштабируется на сложные многообразия.

\section{Многообразие матриц фиксированного ранга}

\subsection{Почему именно ранг $r$, а не «не больше $r$»}

Множество
\[
\mathcal M_r=\set{X\in \R^{m\times n}: \rank X=r}
\]
--- гладкое многообразие. А вот множество
\[
\set{X:\ \rank X\le r}
\]
обычно \emph{не} является гладким многообразием: в точках меньшего ранга меняется локальная размерность, и там возникают сингулярности.

\begin{warningbox}
Для геометрии и оптимизации важно различать:
\begin{itemize}
  \item \textbf{точное ранговое ограничение} $\rank X=r$ --- гладкая страта;
  \item \textbf{нестрогое ограничение} $\rank X\le r$ --- стратифицированное множество с особенностями.
\end{itemize}
\end{warningbox}

\subsection{Размерность}

Если $X=U\Sigma V^T$, где $U\in \St(r,m)$, $V\in \St(r,n)$, $\Sigma\in \R^{r\times r}$ невырождена, то интуитивно:
\begin{itemize}
  \item выбор $U$ даёт $mr-\frac{r(r+1)}{2}$ степеней свободы;
  \item выбор $V$ даёт $nr-\frac{r(r+1)}{2}$;
  \item выбор диагонально-подобного «сердца» ранга $r$ даёт $r^2$.
\end{itemize}
После упрощения получаем
\[
\dim \mathcal M_r = r(m+n-r).
\]

Эту формулу важно помнить.

\subsection{Касательное пространство}

Пусть компактное SVD точки $X$ имеет вид
\[
X=U\Sigma V^T,
\qquad U\in \St(r,m),\ \Sigma\in \R^{r\times r},\ V\in \St(r,n).
\]
Тогда любое касательное направление можно записать как
\[
\Xi = U M V^T + U_\perp B V^T + U C V_\perp^T,
\]
где
\begin{itemize}
  \item $M\in \R^{r\times r}$;
  \item $B\in \R^{(m-r)\times r}$;
  \item $C\in \R^{r\times (n-r)}$.
\end{itemize}

Эквивалентно,
\[
T_X\mathcal M_r = \set{UMV^T + U_pV^T + UV_p^T:\ U^TU_p=0,\ V^TV_p=0}.
\]

Смысл decomposition такой:
\begin{itemize}
  \item член $UMV^T$ меняет сингулярные значения и внутреннюю координату в текущем ранге;
  \item $U_pV^T$ меняет столбцовое пространство;
  \item $UV_p^T$ меняет строковое пространство.
\end{itemize}

\subsection{Проектор на касательное пространство}

Как уже отмечалось, при $P_U=UU^T$ и $P_V=VV^T$
\[
\Proj_{T_X\mathcal M_r}(Z)=P_UZ + ZP_V - P_UZP_V.
\]

Если $\bar f$ --- функция в ambient-пространстве матриц, то риманов градиент на $\mathcal M_r$ при стандартной метрике Фробениуса равен
\[
\grad f(X)=\Proj_{T_X\mathcal M_r}(\nabla \bar f(X)).
\]

\subsection{Ретракции}

Наиболее естественная ретракция:
\[
\Retr_X(\Xi)=\Pi_r(X+\Xi),
\]
где $\Pi_r$ --- усечённое SVD, то есть лучшая в норме Фробениуса матрица ранга $r$, приближающая $X+\Xi$.

Почему это разумно? После малого касательного шага $X+\Xi$ обычно уже не имеет ранг ровно $r$, но находится рядом с многообразием. Усечённое SVD возвращает нас обратно на гладкую страту ранга $r$.

\subsection{Где это используется}

Низкоранговые многообразия возникают в:
\begin{itemize}
  \item matrix completion;
  \item recommendation systems;
  \item reduced-order models;
  \item низкоранговых аппроксимациях больших операторов;
  \item задачах обучения с факторизованными параметрами.
\end{itemize}

Во всех этих задачах риманова геометрия позволяет делать оптимизацию прямо по множеству матриц нужного ранга, а не штрафовать ранг косвенными суррогатами.

\section{Многообразие Stiefel}

\subsection{Определение и размерность}

\begin{definition}
Многообразие Stiefel определяется как
\[
\St(p,n)=\set{X\in \R^{n\times p}: X^TX=I_p}.
\]
Его точки --- это наборы из $p$ ортонормированных столбцов в $\R^n$.
\end{definition}

Ограничение $X^TX=I_p$ задаёт $\frac{p(p+1)}{2}$ независимых скалярных условий, поэтому
\[
\dim \St(p,n)=np-\frac{p(p+1)}{2}.
\]

\subsection{Геометрический смысл}

Точка Stiefel --- это \emph{ортонормированный фрейм}: упорядоченный набор $p$ ортонормированных векторов. Поэтому Stiefel помнит не только само подпространство, но и выбор конкретного базиса в нём.

Если забыть выбор базиса и помнить только натянутое подпространство, мы переходим к Grassmann.

\subsection{Касательное пространство}

Дифференцируем ограничение $X^TX=I$ в направлении $\Xi$:
\[
X^T\Xi + \Xi^T X = 0.
\]
Отсюда
\[
\boxed{\ T_X\St(p,n)=\set{\Xi\in \R^{n\times p}: X^T\Xi + \Xi^TX = 0}.\ }
\]

Это означает, что матрица $X^T\Xi$ должна быть кососимметричной.

\subsection{Проекция и риманов градиент}

При стандартной метрике Фробениуса
\[
\Proj_{T_X\St}(Z)=Z-X\Sym(X^TZ).
\]
Следовательно,
\[
\grad f(X)=\nabla \bar f(X)-X\Sym(X^T\nabla \bar f(X)).
\]

\subsection{Ретракции}

На практике чаще всего используются:
\begin{enumerate}
  \item \textbf{QR-ретракция:}
  \[
  \Retr_X(\Xi)=\operatorname{qf}(X+\Xi);
  \]
  \item \textbf{полярная ретракция:}
  \[
  \Retr_X(\Xi)=(X+\Xi)(I+\Xi^T\Xi)^{-1/2}.
  \]
\end{enumerate}

QR дешевле, полярная часто геометрически аккуратнее.

\subsection{Где Stiefel появляется в приложениях}

\begin{itemize}
  \item задачи на ортогональные матрицы;
  \item orthogonal Procrustes;
  \item constrained PCA и подпространственные методы;
  \item ICA, CCA, некоторые нейросетевые модели с ортогональными весами;
  \item оптимизация по фреймам в численных методах.
\end{itemize}

\section{Grassmann: пространство подпространств}

\subsection{Что такое точка Grassmann}

\begin{definition}
Grassmann-многообразие $\Gr(p,n)$ --- это множество всех $p$-мерных линейных подпространств пространства $\R^n$.
\end{definition}

Если $X\in \St(p,n)$, то его столбцы порождают некоторое подпространство $\operatorname{span}(X)$. При правом умножении на ортогональную матрицу $Q\in O(p)$ подпространство не меняется:
\[
\operatorname{span}(XQ)=\operatorname{span}(X).
\]
Поэтому
\[
\Gr(p,n)=\St(p,n)/O(p).
\]

\subsection{Размерность}

У Stiefel мы помним ещё и выбор ортонормированного базиса, а группа $O(p)$ отвечает за его переобозначение. Так как
\[
\dim O(p)=\frac{p(p-1)}{2},
\]
получаем
\[
\dim \Gr(p,n)=np-\frac{p(p+1)}{2}-\frac{p(p-1)}{2}=p(n-p).
\]

\subsection{Касательное пространство}

Если фиксировать представителя $X\in \St(p,n)$, то касательные направления к Grassmann удобно выбирать в виде матриц, ортогональных столбцовому пространству $X$:
\[
\boxed{\ T_{[X]}\Gr(p,n)\cong \set{\Xi\in \R^{n\times p}: X^T\Xi=0}.\ }
\]

Здесь квадратные движения вида $X\Omega$ с кососимметричной $\Omega$ меняют лишь базис внутри того же подпространства и потому считаются «вертикальными», а не настоящими движениями по Grassmann.

\subsection{Проекция}

При стандартной метрике
\[
\boxed{\ \Proj_{T_{[X]}\Gr}(Z)=(I-XX^T)Z.\ }
\]
Поэтому
\[
\grad f([X])=(I-XX^T)\nabla \bar f(X).
\]

\subsection{Геодезика и ретракции}

Если $\Xi\in T_{[X]}\Gr$ и его SVD равно
\[
\Xi = U\Sigma V^T,
\]
то точная геодезика имеет красивую формулу
\[
\Gamma(t)= X V\cos(t\Sigma)V^T + U\sin(t\Sigma)V^T.
\]
Соответственно,
\[
\Exp_{[X]}(\Xi)=\Gamma(1).
\]

Но на практике почти всегда используют более дешёвые ретракции, например
\[
\Retr_X(\Xi)=\operatorname{qf}(X+\Xi).
\]

\subsection{Где Grassmann особенно естественен}

Grassmann возникает, когда объектом является \emph{не базис, а подпространство}. Например:
\begin{itemize}
  \item PCA и subspace tracking;
  \item invariant subspace problems;
  \item методы на подпространствах Крылова;
  \item задачи типа dominant eigenspace / low-dimensional representation.
\end{itemize}

\section{Как всё это связано в одну картину}

Теперь можно увидеть общую логику.

\subsection{Шаг 1: допустимое множество --- это многообразие}

Сначала мы замечаем, что ограничения задачи задают гладкое множество допустимых точек. Часто оно описывается уравнениями $F(x)=0$ с полным рангом Якоби, а значит образует гладкое подмногообразие в $\R^n$.

\subsection{Шаг 2: локальная линейная аппроксимация --- касательное пространство}

В точке $x$ многообразие локально аппроксимируется пространством $T_xM$. Это и есть пространство всех «мгновенно допустимых» направлений.

\subsection{Шаг 3: чтобы говорить о наискорейшем спуске, нужна метрика}

Один и тот же дифференциал $df_x$ при разных метриках даёт разные градиенты. Поэтому выбор метрики не декоративен: он определяет геометрию движения по многообразию.

\subsection{Шаг 4: embedded-геометрия даёт простую формулу}

Если $M\subset \R^n$ и метрика индуцирована из объёмлющего пространства, то
\[
\grad f(x)=\Proj_{T_xM}(\nabla \bar f(x)).
\]

То есть нам достаточно уметь:
\begin{enumerate}
  \item считать обычный евклидов градиент;
  \item проектировать его на касательное пространство.
\end{enumerate}

\subsection{Шаг 5: нужно вернуть точку обратно на многообразие}

Касательное пространство --- линейная аппроксимация, а многообразие криволинейно. Поэтому после шага вдоль касательного вектора нужен механизм возврата на $M$: экспонента или ретракция.

\subsection{Шаг 6: отсюда и появляется риманов градиентный спуск}

Формула
\[
x_{k+1}=\Retr_{x_k}\!\bigl(-\alpha_k\grad f(x_k)\bigr)
\]
просто собирает все предыдущие идеи в одну вычислительную схему.

\section{Часто встречающиеся примеры в одной таблице}

\renewcommand{\arraystretch}{1.28}
\begin{longtable}{>{\raggedright\arraybackslash}p{0.2\textwidth} >{\raggedright\arraybackslash}p{0.27\textwidth} >{\raggedright\arraybackslash}p{0.24\textwidth} >{\raggedright\arraybackslash}p{0.21\textwidth}}
\toprule
\textbf{Многообразие} & \textbf{Касательное пространство} & \textbf{Проектор} & \textbf{Простая ретракция} \\
\midrule
\endfirsthead
\toprule
\textbf{Многообразие} & \textbf{Касательное пространство} & \textbf{Проектор} & \textbf{Простая ретракция} \\
\midrule
\endhead
$S^{n-1}$ & $\set{\xi:\ x^T\xi=0}$ & $z\mapsto z-(x^Tz)x$ & $\dfrac{x+\xi}{\norm{x+\xi}}$ \\
\midrule
$\St(p,n)$ & $\set{\Xi:\ X^T\Xi+\Xi^TX=0}$ & $Z\mapsto Z-X\Sym(X^TZ)$ & $\operatorname{qf}(X+\Xi)$ \\
\midrule
$\Gr(p,n)$ & $\set{\Xi:\ X^T\Xi=0}$ & $Z\mapsto (I-XX^T)Z$ & $\operatorname{qf}(X+\Xi)$ \\
\midrule
$\mathcal M_r$ & $UMV^T+U_pV^T+UV_p^T$ & $Z\mapsto P_UZ+ZP_V-P_UZP_V$ & $\Pi_r(X+\Xi)$ \\
\bottomrule
\end{longtable}

\section{На что обратить внимание при повторении темы}

\begin{enumerate}
  \item \textbf{Не путать пространство и его локальную модель.}
  Многообразие локально похоже на $\R^m$, но обычно не является линейным пространством.

  \item \textbf{Не путать дифференциал и градиент.}
  $df_x$ --- ковектор, $\grad f(x)$ --- вектор, зависящий от метрики.

  \item \textbf{Для embedded-задач помнить золотую формулу}
  \[
  \grad f(x)=\Proj_{T_xM}(\nabla \bar f(x)).
  \]

  \item \textbf{Для множеств, заданных уравнениями, сначала искать касательное пространство.}
  Часто именно это превращает геометрическую задачу в линейную алгебру.

  \item \textbf{Ретракция нужна не вместо геометрии, а как её вычислительная аппроксимация.}

  \item \textbf{Grassmann и Stiefel не одно и то же.}
  Stiefel хранит \emph{ортонормированный базис}, Grassmann --- только \emph{подпространство}.

  \item \textbf{Множество точного ранга $r$ гладкое, а множество ранга не больше $r$ --- уже нет.}
\end{enumerate}

\section{Короткая шпаргалка формул}

\begin{summarybox}
\textbf{Подмногообразие по ограничениям:}
\[
M=\set{x:F(x)=0},\qquad T_xM=\ker dF_x.
\]

\textbf{Риманов градиент:}
\[
df_x[\xi]=g_x(\grad f(x),\xi).
\]

\textbf{В координатах:}
\[
(\grad f)^i=\sum_j g^{ij}\partial_j f.
\]

\textbf{Для индуцированной метрики:}
\[
\grad f(x)=\Proj_{T_xM}(\nabla \bar f(x)).
\]

\textbf{Риманов шаг спуска:}
\[
x_{k+1}=\Retr_{x_k}\!\bigl(-\alpha_k\grad f(x_k)\bigr).
\]

\textbf{Stiefel:}
\[
T_X\St=\set{\Xi:X^T\Xi+\Xi^TX=0},
\qquad
\Proj(Z)=Z-X\Sym(X^TZ).
\]

\textbf{Grassmann:}
\[
T_{[X]}\Gr\cong\set{\Xi:X^T\Xi=0},
\qquad
\Proj(Z)=(I-XX^T)Z.
\]

\textbf{Fixed-rank:}
\[
\Proj(Z)=P_UZ+ZP_V-P_UZP_V.
\]
\end{summarybox}

\section*{Послесловие}
\addcontentsline{toc}{section}{Послесловие}

Если смотреть на всю конструкцию целиком, то риманова оптимизация оказывается очень естественной.

Гладкое многообразие даёт язык для описания допустимого множества. Касательное пространство описывает допустимые мгновенные направления движения. Риманова метрика позволяет измерять длины и определять наискорейший спуск. Проекция переносит ambient-градиент в геометрию допустимых направлений. Экспонента и ретракция возвращают нас с линейной касательной модели обратно на само многообразие. А матричные многообразия --- Stiefel, Grassmann и фиксированный ранг --- превращают эту схему в реальный вычислительный инструмент для современных задач линейной алгебры и оптимизации.

Именно поэтому полезно держать в голове не набор разрозненных формул, а следующую цепочку:
\[
\text{ограничение} \ \Longrightarrow \ \text{многообразие} \ \Longrightarrow \ T_xM \ \Longrightarrow \ g_x \ \Longrightarrow \grad f \ \Longrightarrow \Retr.
\]

Когда эта логика становится прозрачной, большинство формул уже не приходится заучивать отдельно: они начинают вытекать из общей картины.

\end{document}
